\documentclass[11pt]{article}

\usepackage[T1]{fontenc}
\usepackage[utf8]{inputenc}
\usepackage{lmodern}
\usepackage{amsmath,amssymb,amsthm}
\usepackage{graphicx}
\usepackage{booktabs}
\usepackage[numbers,sort&compress]{natbib}
\usepackage[colorlinks=true,linkcolor=blue,citecolor=blue,urlcolor=blue]{hyperref}

\title{TITLE PLACEHOLDER}
\author{Author Name}
\date{\today}

\begin{document}

\maketitle

\begin{abstract}
Summarize the problem, the approach, and the main result in a few sentences.
\end{abstract}

\section{Introduction}\label{sec:intro}

Introduce the problem and state your contributions. Prior work can be cited
with \citep{sample2026} for parenthetical citations or \citet{sample2026} for
textual ones.

\section{Methods}\label{sec:methods}

Describe the approach. Equations go in the \texttt{equation} environment:

\begin{equation}\label{eq:main}
  f(x) = \int_{-\infty}^{\infty} e^{-t^2}\,\mathrm{d}t
\end{equation}

Equation~\eqref{eq:main} is referenced with \verb|\eqref|.

\section{Results}\label{sec:results}

Tables use \texttt{booktabs} rules:

\begin{table}[htbp]
  \centering
  \begin{tabular}{lcc}
    \toprule
    Method & Metric A & Metric B \\
    \midrule
    Baseline & 0.50 & 0.61 \\
    Ours     & \textbf{0.87} & \textbf{0.92} \\
    \bottomrule
  \end{tabular}
  \caption{Comparison against the baseline.}
  \label{tab:comparison}
\end{table}

Figures go in \texttt{figures/} and are included like this:

\begin{figure}[htbp]
  \centering
  % \includegraphics[width=\linewidth]{figures/example.pdf}
  \caption{An example figure (uncomment \texttt{includegraphics} once added).}
  \label{fig:example}
\end{figure}

\section{Conclusion}\label{sec:conclusion}

Restate the contribution and outline future work.

\bibliographystyle{plainnat}
\bibliography{refs}

\end{document}
